\section*{Abstract}
\addcontentsline{toc}{section}{Abstract}

The alignment of a student's academic capabilities with an appropriate career trajectory is a foundational challenge in modern educational systems. In Rwanda, students transition from Ordinary Level (Senior 3) to Advanced Level based on their national examination performance, a critical juncture that heavily dictates their future career options, particularly within the highly demanding health sector. Traditional guidance systems rely heavily on manual counseling, which is often subjective, constrained by cognitive biases, and incapable of processing large-scale historical performance data. Furthermore, early algorithmic approaches to career guidance typically frame the problem as a mutually exclusive, single-label classification task, functioning as "black boxes" that fail to explain the mathematical reasoning behind their recommendations. 

This thesis develops and evaluates an Explainable Multi-Label Machine Learning Model designed to provide transparent, multi-pathway educational and career guidance. Utilizing a comprehensive longitudinal dataset of Senior 3 national examination results spanning from 2017 to 2025, the study focuses on core STEM subjects (Biology, Chemistry, Mathematics, and Physics). The methodology employs the Classifier Chains technique to account for the inherent correlations between related health sector careers (e.g., Medicine and Nursing). Four distinct base algorithms—Decision Trees, Random Forests, Extreme Gradient Boosting (XGBoost), and Multi-Layer Perceptrons (MLP)—are evaluated based on Subset Accuracy, Hamming Loss, and Micro F1-Scores. 

Crucially, the system moves beyond mere prediction by integrating an Explainable AI (XAI) layer. It utilizes feature importance mapping to provide global interpretability and introduces the mathematical framework for Counterfactual Explanations to offer local, actionable advice (e.g., identifying the minimum required improvement in specific subjects to change a rejected recommendation to an accepted one). The results demonstrate that ensemble methods, combined with the Classifier Chains mechanism, significantly outperform independent binary classifiers in capturing complex, non-linear relationships in student data while preserving high interpretability. This research establishes a robust, mathematically sound, and scalable framework that not only predicts but explicitly explains career recommendations, laying the theoretical groundwork for future integration with longitudinal tracking, continuous adaptation, and dynamic labor market data.

\vspace{1cm}
\textbf{Keywords:} Educational Data Mining, Multi-Label Classification, Classifier Chains, Explainable AI, XGBoost, Random Forest, Counterfactual Explanations, Career Guidance.
